% =====================================================================
% Option A / gates G0+G0.5: EveNet data-efficiency test for kappa_lambda
% density-ratio learning in HH->bbtautau (CMS full simulation).
% Conventions: British -ise spellings, first-occurrence expansion.
% Build: pdflatex twice. No shell-escape, no external assets.
% =====================================================================
\documentclass[11pt]{article}

\usepackage[margin=2.5cm]{geometry}
\usepackage{amsmath,amssymb}
\usepackage{booktabs}
\usepackage{enumitem}
\usepackage{xcolor}
\usepackage{microtype}
\emergencystretch=1.5em
\usepackage[colorlinks=true,linkcolor=blue!60!black,citecolor=blue!60!black,urlcolor=blue!60!black]{hyperref}

\newcommand{\kl}{\ensuremath{\kappa_\lambda}}
\newcommand{\bbtt}{\ensuremath{b\bar b\,\tau^+\tau^-}}
\newcommand{\tauh}{\ensuremath{\tau_h}}
\newcommand{\mhh}{\ensuremath{m_{HH}}}
\newcommand{\IC}{\ensuremath{\mathrm{IC}}}
\newcommand{\SCr}{\ensuremath{\mathrm{SC_{rms}}}}
\newcommand{\SCn}{\ensuremath{\mathrm{SC_{norm}}}}
\newcommand{\Eref}{\ensuremath{\mathbb{E}_{\mathrm{ref}}}}

\title{EveNet for \kl\ Density-Ratio Learning in $HH\to\bbtt$:\\
Pre-training Improves Discrimination\\
but Degrades Low-Statistics Ratio Calibration}
\author{Working note --- \bbtt\ \kl\ NSBI project}
\date{16 July 2026 (rev.\ 14 August: (P5), Sec.~\ref{sec:power};
20 August 2026: verdict, Sec.~\ref{sec:p5}) \\[2pt]
\small Gates G0 / G0.5 of \texttt{EVENET\_FEASIBILITY\_NOTE}; 90-run sweep, branch \texttt{fm}}

\begin{document}
\maketitle

\begin{abstract}
\noindent
We test whether the pre-trained event-level foundation model (FM)
EveNet~\cite{Hsu:2026evenet} improves the \emph{data efficiency} of density-ratio learning
for a Higgs self-coupling (\kl) measurement by neural simulation-based inference (NSBI),
on CMS full-simulation $HH\to\bbtt$ events. Three arms --- fine-tuned, frozen-backbone and
from-scratch --- identical in data, architecture and head and differing \emph{only} in
initialisation, were trained at six training fractions $\times$ five seeds (90 runs) and
evaluated on one fixed held-out split against two axes fixed before unblinding:
discrimination (weighted area under the receiver-operating characteristic curve, AUC) and
ratio calibration (integral and shape closure). Three results. \textbf{(i)~Discrimination:
pre-training wins decisively} --- $+0.014$ to $+0.017$ AUC at every fraction $\ge 1\%$, worth
${\approx}3$--$4.5\times$ the training data, with both pre-trained arms exceeding a
feature-level gradient-boosted decision tree (BDT) ceiling. \textbf{(ii)~Calibration
quality: pre-training wins at full statistics} --- with five-model ensembling the
frozen-backbone arm is the \emph{only} configuration in the study that passes both analysis
gates, and both pre-trained arms show ${\approx}3$--$4\sigma$ less residual structure than
from-scratch. \textbf{(iii)~Calibration data efficiency: resolved --- no improvement, and measurable harm
for the frozen arm} --- the (P5) $K{=}20$ campaign (225 further trainings, paired
test-set bootstrap; Sec.~\ref{sec:p5}) gives
$\Delta_{\rm finetune} = -0.022 \pm 0.026$ (bounded: $|\Delta| \lesssim 0.05$) and
$\Delta_{\rm frozen} = -0.079 \pm 0.031$ ($z = -2.6$, crossing the pre-registered harm
threshold), with the secondary shape statistic agreeing ($z = -2.0$ and $-3.2$). On the
same ensembles the AUC advantage persists, so pre-training buys event ordering while
\emph{costing} low-$N$ calibration --- consistent with an initialisation bias in the
weakly-pinned normalisation direction (Sec.~\ref{sec:power}) that ensembling cannot
remove and that full statistics eventually repins: at fraction 1.0 the frozen ensemble
remains the study's only gate pass. In
both (i) and (ii) the frozen arm matches or beats fine-tuning, locating the benefit in the
pre-trained \emph{representation} rather than in backbone adaptation. Section~\ref{sec:corr}
documents three conversion errors found and fixed during analysis, one of which had
inflated all closure numbers by a factor of ${\approx}3$; a fourth hypothesis, that
EveNet's diffusion-noised objective causes miscalibration, was refuted by source audit
before the ablation was believed.
\end{abstract}

\tableofcontents

% =====================================================================
\section{Setup}
\label{sec:setup}

\paragraph{Task.} Learn per-event density ratios $r(x;\kl,\kl')$ for an unbinned NSBI
likelihood. The \kl\ dependence is a \emph{continuous} deformation of one process's
spectrum (the \mhh\ lineshape and triangle--box interference), not the separation of two
distinct processes.

\paragraph{Data.} CMS 2024 full-simulation $HH\to\bbtt$ Powheg samples through the KIT
CROWN chain (published $\mu\tau_h$/$e\tau_h$ selection), converted to EveNet's stock
7-feature token schema: four tokens per event ($b_1$, $b_2$, \tauh, $\ell$), missing
transverse energy in the global conditioning block, weights $=$ \texttt{genWeight}
$\times$ \texttt{puweight} (signed; the reference carries 6.1\% negative weights, the
hypothesis 1.0\%). Requiring a valid $H\to b\bar b$ candidate drops 18.6--24.5\% of events
(\kl-dependent). Class~0 $=$ \kl$=1$ (reference, 488k), class~1 $=$ \kl$=5$ (339k); merged,
shuffled, split $80/10/10$ into 661.9k train / 82.7k validation / 82.7k test. Training
fractions multiply the \emph{train} split; \textbf{all} runs share the one fixed test split.

\paragraph{Arms.} \textbf{finetune} (20M pre-trained checkpoint, all weights trainable);
\textbf{frozen} (same checkpoint, backbone GlobalEmbedding/PET/ObjectEncoder frozen, heads
trainable --- isolates the representation from optimiser effects); \textbf{scratch} (random
init). Fractions $\{0.003, 0.01, 0.02, 0.05, 0.1, 1.0\}$ ($\approx 2.0$k to 662k training
events) $\times$ 5 seeds $=$ 90 runs. Options files are the authors' own downstream recipes,
vendored unchanged apart from the freeze flags.

\paragraph{External bar.} A gradient-boosted decision tree on the 10 high-level features
behind the published unbinned result (Sec.~\ref{sec:ceiling}).

% =====================================================================
\section{Pre-registration}
\label{sec:prereg}

Fixed before any sweep output existed, reproduced unchanged:

\begin{enumerate}[label=(P\arabic*),leftmargin=2.6em]
\item \textbf{Decision primacy.} The adoption metric is the \emph{closure left-shift}: the
  smallest fraction at which an arm passes both gates, compared between finetune and
  scratch; their ratio is the equivalent-data multiplier. A left-shift is an adoption-grade
  win \emph{even at AUC parity}; an AUC-only win with failing closure is \emph{not} a win.
  AUC is a secondary no-regression guard.
\item \textbf{Gates.} $|\IC| \le \max(1\%,\,3\sigma_{\mathrm{stat}})$ \emph{and}
  $\SCr \le 0.05$ ($=$ the analysis \texttt{CLOSURE\_TOL}); $\chi^2/\mathrm{ndf}\approx 1$
  indicates noise-only residuals.
\item \textbf{Checkpoint selection.} Every arm is evaluated at its best-validation-loss
  checkpoint (EveNet's predictor otherwise loads the final, most overfit, epoch).
\item \textbf{Fallback.} If no arm passes anywhere, leaving the left-shift undefined, the
  secondary reading is relative closure at matched fractions, reported as such.
\item \textbf{Seed-extension decision rule} \emph{(added 14 August 2026; fixed before any
  $K>5$ output existed).} The sweep is extended by seeds $5$--$19$ at fractions
  $\{0.003, 0.01, 0.02, 0.05, 0.1\}$ --- 225 runs, identical data, recipe,
  best-validation selection and test split --- giving $K{=}20$ ensembles per cell there.
  For each pre-trained arm $a \in \{\text{finetune}, \text{frozen}\}$ define
  \[ \Delta_a = \tfrac{1}{5}\sum_f \big( |\IC|^{\mathrm{scratch}}_f - |\IC|^{a}_f \big),
     \qquad
     \sigma_{\Delta}^2 = \tfrac{1}{25}\sum_f \Big[\,
     \sigma^2_{\mathrm{boot},f} \;+\;
     \tfrac{\hat\sigma^2_{a,f} + \hat\sigma^2_{\mathrm{scratch},f}}{20} \,\Big], \]
  where $|\IC|_f$ are the $K{=}20$ ensemble values, $\sigma_{\mathrm{boot},f}$ is the
  standard deviation over 200 Poisson($\mu{=}1$) test-set replicas of the \emph{paired}
  arm difference (common draws across arms at each fraction), and $\hat\sigma$ are the
  per-seed spreads. Verdicts: $\Delta_a > 2\sigma_\Delta$ --- pre-training improves ratio
  calibration; $\Delta_a < -2\sigma_\Delta$ --- it harms it; otherwise any effect is
  bounded and quoted as $|\Delta_a| \lesssim 2\sigma_\Delta$ (forecast ${\approx}0.05$).
  The same statistic on \SCr\ is reported as secondary, not decisive. All three outcomes
  are reportable; none will be reinterpreted after unblinding.
  \emph{[Executed 20 August 2026; verdict in Sec.~\ref{sec:p5}.]}
\end{enumerate}

\noindent (P4) is the branch that materialised for the finetune--scratch comparison. (P5)
was registered after the $K{=}5$ results below were known, but before any extension job
was submitted; the $K{=}5$ data motivated the campaign and play no further role in the
verdict.

% =====================================================================
\section{Metrics}
\label{sec:metrics}

\subsection{The closure identity}

The ratio estimate is $\hat r = \frac{p}{1-p}\cdot f_{\mathrm{prior}}$, with
$f_{\mathrm{prior}}$ fixed by the convention EveNet's loss implies
(Sec.~\ref{sec:corr}). Everything rests on one identity: averaging the \emph{true} ratio
over reference events,
%
\begin{equation}
\Eref[r] = \int p_{\mathrm{ref}}(x)\,\frac{p_{\mathrm{hyp}}(x)}{p_{\mathrm{ref}}(x)}\,
\mathrm{d}x = \int p_{\mathrm{hyp}}(x)\,\mathrm{d}x = 1 .
\label{eq:closure}
\end{equation}
%
The reference density cancels and what remains integrates to one. Intuitively $r(x)$ says
``this event is $r$ times more likely under the hypothesis''; over the reference ensemble
the $r>1$ and $r<1$ regions must balance exactly. On weighted Monte Carlo the expectation
becomes the self-normalised weighted average
$\sum_{i\in\mathrm{ref}} w_i \hat r_i / \sum_{i\in\mathrm{ref}} w_i$, converging at
$1/\sqrt{n_{\mathrm{eff}}}$ --- so closure is always ``consistent with 1'', never ``equal
to 1''.

\subsection{The two metrics}

\begin{description}[leftmargin=2.4em,style=nextline]
\item[{Integral closure, $\IC = \Eref[\hat r]-1$}]
  Reweight the reference by $\hat r$ and ask whether the \emph{total} hypothesis yield is
  recovered. $\IC = +0.20$ means a 20\% over-production. This is the \emph{normalisation}
  error. Quoted with its weighted standard error.
\item[{Shape closure, \SCr\ and \SCn}]
  $\IC$ can be satisfied by errors that cancel globally, so Eq.~\eqref{eq:closure} is
  applied \emph{differentially}: quantile-bin the score, and in each bin compare the
  $\hat r$-reweighted reference mass $A_b$ with the hypothesis mass $B_b$, forming
  $\mathrm{rel}_b = A_b/B_b - 1$. We report the RMS over bins with each bin's statistical
  variance subtracted in quadrature. We bin in the classifier score because it is monotone
  in $\hat r$, so miscalibration is forced to appear there.
\end{description}

\paragraph{Why the per-bin error is $1/n_{\mathrm{eff}}$.} For a weighted sum
$S=\sum_i v_i$ with independent (Poisson) presence, $\mathrm{Var}(S)=\sum_i v_i^2$, so
%
\begin{equation}
\frac{\mathrm{Var}(S)}{S^2}=\frac{\sum_i v_i^2}{(\sum_i v_i)^2}\equiv\frac{1}{n_{\mathrm{eff}}},
\qquad n_{\mathrm{eff}}\equiv\frac{(\sum_i v_i)^2}{\sum_i v_i^2},
\end{equation}
%
which \emph{defines} the effective sample size (Kish). Unweighted events give
$n_{\mathrm{eff}}=N$ and the Poisson $1/\sqrt N$; unequal weights always reduce it
($v=(10,1,1,1)$ gives $n_{\mathrm{eff}}=1.6$, four events behaving like 1.6). On the
reference side $v_i = w_i\hat r_i$, so the reweighting inflates the spread --- which is why
bins at $\hat r\to\infty$ are dropped by the $n_{\mathrm{eff}}$ threshold rather than
allowed to dominate. Since $\mathrm{rel}_b$ is a ratio of estimates from two
\emph{independent} samples, $\sigma_b^2 = 1/n_{\mathrm{eff},0}+1/n_{\mathrm{eff},1}$.

\paragraph{Why subtract in quadrature.} With
$\mathrm{rel}_b=\beta_b+\varepsilon_b$ ($\beta$ bias, $\varepsilon$ noise of variance
$\sigma_b^2$), $\mathbb{E}[\mathrm{rel}_b^2]=\beta_b^2+\sigma_b^2$: the naive RMS is
inflated by noise, worst at small statistics. Subtracting the known $\sigma_b^2$ estimates
the bias RMS, $\SCr=\sqrt{\max(\langle\mathrm{rel}^2-\sigma^2\rangle,0)}$. The clipping at
zero is why exact $0.000$ entries appear: they mean ``no bias resolvable above the
statistical floor''. Without the subtraction the metric rejects correct models --- the
first implementation did exactly that, failing an analytically calibrated toy in three
seeds of four. The same $\sigma_b$ normalises
$\chi^2/\mathrm{ndf}=\langle(\mathrm{rel}_b/\sigma_b)^2\rangle$.

\subsection{\SCr\ versus \SCn}

\SCr\ is computed from $\hat r$ as delivered, \SCn\ from $\hat r/\Eref[\hat r]$. For a
\emph{pure} scale error $\hat r = c\,r_{\mathrm{true}}$, every bin has $A_b/B_b = c$, so
$\SCr\simeq|c-1|\simeq|\IC|$ although the shape is perfect, while $\SCn = 0$. The two raw
numbers are therefore not independent, and reading them together gives three signatures:

\begin{table}[h]
\centering\small
\begin{tabular}{@{}lcccl@{}}
\toprule
signature & $|\IC|$ & \SCr & \SCn & diagnosis \\
\midrule
calibrated & $\approx 0$ & $\approx 0$ & $\approx 0$ & usable \\
pure scale error & $|c-1|$ & $\approx|\IC|$ & $\approx 0$ & normalisation only --- removable \\
shape distortion & any & large & large & deformed likelihood --- not removable \\
\bottomrule
\end{tabular}
\caption{Validated on synthetic data: a $\times1.4$ scale error gives
$(0.399,\,0.399,\,0.000)$; a temperature-$1.5$ distortion gives $(0.450,\,0.694,\,0.503)$.}
\label{tab:signatures}
\end{table}

\paragraph{Why the distinction is physical.} Process $p$ enters the likelihood as
$\nu_p\,p_p(x)$ with $\nu_p$ the yield (from cross section $\times$ luminosity $\times$
acceptance, determined independently of the classifier) and $p_p = r_p\,p_{\mathrm{ref}}$.
If $\hat r_p = c\,r_p$ then $\nu_p\hat p_p = (c\nu_p)[\hat p_p/c]$: the constant is
algebraically indistinguishable from mis-stating the \emph{yield}, leaving the normalised
shape untouched. It is therefore removable by construction (rescale by $\Eref[\hat r]$ on
held-out events) and, uncorrected, propagates only through that process's normalisation
where it is degenerate with a rate nuisance parameter. For a coupling measurement this is
not harmless --- $\sigma(\kl)$ carries the physics --- but it is a one-parameter defect with
a one-parameter cure. A shape error is different in kind: $\hat r(x)/r(x)$ varies with $x$,
events are mis-weighted relative to one another, and no constant or rate parameter repairs
it. \SCr\ is the pre-registered gate; \SCn\ is a post-hoc decomposition reported alongside.

\paragraph{What closure cannot catch.} Both metrics test \emph{self-consistency}, not
\emph{sufficiency}. A classifier on lossy features $\phi(x)$ converges to the exact ratio of
the \emph{marginalised} densities, satisfies Eq.~\eqref{eq:closure} identically, and passes
any test binned in a function of $\phi$ --- the degenerate $\hat r\equiv 1$ passes both
metrics perfectly while carrying no information. Information loss appears not as bias but as
a \emph{flatter} ratio, i.e.\ lost discrimination. Closure catches lies, not omissions: it
earns $\hat r$ the right to enter the likelihood, while AUC shows it is worth entering.

% =====================================================================
\section{Result 1 --- discrimination (G0)}
\label{sec:auc}

\begin{table}[h]
\centering\small
\begin{tabular}{@{}lccc|ccc@{}}
\toprule
& \multicolumn{3}{c|}{\textbf{single model} (seed mean $\pm$ s.d.)}
& \multicolumn{3}{c}{\textbf{5-seed ensemble}} \\
fraction & finetune & frozen & scratch & finetune & frozen & scratch \\
\midrule
0.003 & 0.7427(57) & 0.7498(80) & \textbf{0.7573(122)} & 0.7553 & 0.7674 & \textbf{0.7698} \\
0.01  & 0.7768(108) & \textbf{0.7793(76)} & 0.7612(300) & \textbf{0.7924} & 0.7914 & 0.7787 \\
0.02  & \textbf{0.7941(43)} & 0.7931(80) & 0.7804(19) & 0.8040 & \textbf{0.8041} & 0.7845 \\
0.05  & 0.8056(54) & \textbf{0.8099(22)} & 0.7949(8) & 0.8144 & \textbf{0.8160} & 0.7983 \\
0.1   & 0.8163(22) & \textbf{0.8169(9)} & 0.8003(25) & 0.8210 & \textbf{0.8212} & 0.8059 \\
1.0   & \textbf{0.8313(2)} & \textbf{0.8313(5)} & 0.8250(4) & \textbf{0.8332} & 0.8330 & 0.8285 \\
\bottomrule
\end{tabular}
\caption{Weighted AUC. Parenthesised digits are the seed standard deviation in units of the
last decimals.}
\label{tab:auc}
\end{table}

\begin{enumerate}[leftmargin=2.2em]
\item \textbf{Pre-training wins decisively.} At fractions $0.01$--$1.0$ both pre-trained
  arms exceed scratch by $0.014$--$0.017$ AUC against seed standard deviations of
  $0.001$--$0.008$ --- many standard deviations, monotone across the grid.
\item \textbf{Worth ${\approx}3$--$4.5\times$ the data.} Interpolating scratch's ensemble
  curve logarithmically, finetune's AUC at $0.01$, $0.02$, $0.05$, $0.1$ requires scratch
  fractions of $0.033$, $0.079$, $0.222$, $0.437$ --- multipliers $3.3$, $4.0$, $4.4$,
  $4.4\times$; at full statistics scratch never reaches finetune's value.
\item \textbf{Frozen $\approx$ finetune everywhere} (differences $\le 0.004$). The benefit
  lives in the pre-trained \emph{representation}; adapting the backbone adds nothing
  measurable. This answers the ``did the pretext retain \kl-relevant continuous
  information?'' question: it did.
\end{enumerate}

\paragraph{The 0.003 inversion.} At the smallest fraction scratch leads. That tier is an
optimisation regime, not a representational one: with ${\approx}2$k events and a global
batch of ${\approx}8$k across 4 GPUs, one ``epoch'' is a single gradient step, and the
selected checkpoints sit at epochs 4--21, i.e.\ the models are compared after
$\mathcal{O}(10)$ updates under \emph{different} learning-rate recipes. We treat it as a
recipe artefact pending a batch-size/learning-rate ablation.

% =====================================================================
\section{Result 2 --- calibration (G0.5)}
\label{sec:closure}

All numbers below use the ratio convention derived in Sec.~\ref{sec:corr}. Production NSBI
uses ensembles (the analysis config trains five members) and the BDT ceiling is itself an
ensemble, so the \textbf{ensemble row is the like-for-like object}; single-model numbers are
given for reference.

\begin{table}[h]
\centering\small
\begin{tabular}{@{}lccc|ccc@{}}
\toprule
& \multicolumn{3}{c|}{single model, $|\IC|$ (seed mean)}
& \multicolumn{3}{c}{single model, \SCn} \\
fraction & finetune & frozen & scratch & finetune & frozen & scratch \\
\midrule
0.003 & 0.136 & 0.310 & 0.164 & 0.328 & 0.476 & 0.365 \\
0.01  & 0.110 & 0.279 & 0.213 & 0.346 & 0.341 & \textbf{0.206} \\
0.02  & 0.148 & 0.120 & 0.129 & 0.276 & 0.244 & \textbf{0.187} \\
0.05  & 0.102 & 0.099 & 0.158 & 0.178 & 0.206 & \textbf{0.159} \\
0.1   & 0.098 & 0.143 & 0.192 & 0.176 & \textbf{0.168} & 0.180 \\
1.0   & 0.050 & 0.053 & 0.055 & 0.090 & \textbf{0.061} & 0.076 \\
\bottomrule
\end{tabular}
\caption{Single-model closure. Statistical errors on $|\IC|$ are $0.009$--$0.021$. No single
model passes the gates at any fraction.}
\label{tab:closure}
\end{table}

\begin{table}[h]
\centering\small
\begin{tabular}{@{}lcccccc@{}}
\toprule
arm & fraction & $|\IC|$ & gate & \SCr & gate & verdict \\
\midrule
finetune & 1.0 & 0.0198 & 0.0559 \checkmark & 0.0521 & 0.05 $\times$ & fails by 4\% \\
\textbf{frozen} & \textbf{1.0} & \textbf{0.0527} & \textbf{0.0536} \checkmark
  & \textbf{0.0272} & \textbf{0.05} \checkmark & \textbf{PASSES} \\
scratch & 1.0 & 0.0362 & 0.0516 \checkmark & 0.0742 & 0.05 $\times$ & fails \\
\bottomrule
\end{tabular}
\caption{Ensemble closure at full statistics against the pre-registered gates
($|\IC|\le\max(1\%,3\sigma_{\mathrm{stat}})$, $\SCr\le0.05$). No arm passes at any smaller
fraction. $\chi^2/\mathrm{ndf}$: finetune $1.80$, frozen $1.52$, scratch $3.20$.}
\label{tab:gates}
\end{table}

\paragraph{(a) Ensembling works --- once the convention is right.} Because $\Eref$ is linear
in $r$, $\IC_{\mathrm{ens}}=\frac1K\sum_k \IC_k$ identically: the ensemble reduces $|\IC|$
only if the members' offsets differ in sign. Under the superseded convention 15 of 18 cells
were same-sign and ensembling did nothing. With the corrected convention the residuals
straddle zero (several signed means are negative) and the ensemble collapses them:
finetune at full statistics goes $0.050\to\mathbf{0.0198}$ (statistical error $0.0186$,
i.e.\ consistent with zero), scratch at $0.01$ goes $0.213\to0.068$.

\paragraph{(b) At full statistics the pre-trained arms are better calibrated in shape.}
Ordering at fraction 1.0 is consistent across all three shape measures: frozen $<$ finetune
$<$ scratch on \SCr\ ($0.027/0.052/0.074$) and on $\chi^2/\mathrm{ndf}$
($1.52/1.80/3.20$). With 20 bins, $\sigma(\chi^2/\mathrm{ndf})\approx\sqrt{2/20}=0.32$, so
scratch's residuals are $\mathbf{3\text{--}4\sigma}$ more structured than either pre-trained
arm's. \textbf{The frozen ensemble is the only configuration in the study that passes both
gates} (Table~\ref{tab:gates}).

\paragraph{(c) But there is no data-efficiency gain --- the pre-registered metric is a null.}
(P1) compares finetune against scratch: neither passes at \emph{any} fraction, so both
first-passing fractions are undefined and no equivalent-data multiplier exists. Under the
(P4) fallback the reduced-fraction ordering is mixed --- scratch is nominally best on \SCn\
at $0.01$--$0.05$, the pre-trained arms at $0.1$--$1.0$, and frozen is the \emph{worst} arm
at $0.003$--$0.01$ while being the best at $1.0$. There is no left-shift anywhere.

\paragraph{Summary.} A usable density ratio requires the full 662k training events
\emph{and} ensembling. At that point pre-training measurably improves ratio shape; it does
not reduce the data needed to get there.

% =====================================================================
\section{The ceiling comparison}
\label{sec:ceiling}

The feature-level BDT, measured at full statistics with class-normalised weights (so its
prior factor is unity by construction):

\begin{table}[h]
\centering\small
\begin{tabular}{@{}lccc@{}}
\toprule
fraction & AUC & $|\IC|$ & shape \\
\midrule
0.01 & 0.7917 & 0.202 & 0.391 \\
0.1  & 0.8130 & 0.011 & 0.039 \\
1.0  & 0.8174 & 0.018 & 0.000 \\
\bottomrule
\end{tabular}
\caption{Kinematic ceiling (10 published features, gradient-boosted trees).}
\label{tab:ceilingtab}
\end{table}

\textbf{Discrimination: EveNet wins.} At $0.1$ and $1.0$ both pre-trained arms exceed the
ceiling ($0.821$/$0.821$ vs.\ $0.813$; $0.833$ vs.\ $0.817$), i.e.\ the object-level model
extracts \kl\ information the 10 aggregate features discard. \emph{Caveat:} the ceiling was
measured on the un-cut flat-feature population whereas the sweep applies the valid-$b$-pair
requirement, so the sweep population is slightly easier; the arm-to-arm comparisons share
identical events and carry no such caveat.

\textbf{Calibration: comparable, once conventions match.} At full statistics the ceiling
reaches $|\IC|=0.018$ and shape $0.000$; the EveNet ensembles reach $|\IC|=0.020$--$0.053$
and $\SCr=0.027$--$0.074$. The frozen ensemble ($0.053$, $0.027$) is within a factor of a
few of the BDT on both axes. An earlier draft of this note reported the BDT as
${\approx}8\times$ better calibrated \emph{despite worse AUC} and drew a strong conclusion
from it; that contrast was an artefact of converting the two estimators under different
conventions (Sec.~\ref{sec:corr}) and does not survive.

% =====================================================================
\section{Why calibration is harder to power than discrimination}
\label{sec:power}

The split verdict of Secs.~\ref{sec:auc}--\ref{sec:closure} --- discrimination resolved
decisively by the same 90 runs that leave the calibration ordering ambiguous --- is not an
accident of this dataset. The two metric families admit structurally different amounts of
noise, on both the training and the test side. This section records the mechanisms and the
measurements behind that statement: they fix the design of (P5), and they generalise to
any foundation-model claim about NSBI.

\paragraph{(a) AUC is a rank statistic; \IC\ is exponentially sensitive to the coordinate
AUC discards.} AUC is invariant under every monotone remapping of the score: it sees only
the ordering of test events. Decompose the difference between two seeds' learnt logit
functions into a monotone part --- a global offset $\delta$ and gain --- and an
order-changing part. The monotone part moves \IC\ directly
($r \propto e^{f}$, so $\IC \approx e^{\delta}-1 \approx \delta$) and moves AUC by exactly
zero; the order-changing part, which is all AUC can see, is small, because two networks
trained on the same data mostly disagree about \emph{how confident} to be, not about which
event is more signal-like. Measured on the sweep's own 90 networks
(Table~\ref{tab:seednoise}), the per-seed spread of \IC\ is $13$--$160\times$ that of AUC,
the disparity \emph{growing} with training statistics.

\begin{table}[h]\centering\small
\begin{tabular}{@{}lcccc@{}}
\toprule
fraction & seed s.d.\ of \IC\ (ft/fr/scr) & seed s.d.\ of AUC (ft/fr/scr) & ratio \\
\midrule
0.003 & 0.145 / 0.219 / 0.247 & 0.006 / 0.008 / 0.012 & $24\times$ \\
0.01  & 0.098 / 0.252 / 0.256 & 0.011 / 0.008 / 0.030 & $13\times$ \\
0.02  & 0.104 / 0.131 / 0.131 & 0.004 / 0.008 / 0.002 & $26\times$ \\
0.05  & 0.094 / 0.081 / 0.152 & 0.005 / 0.002 / 0.001 & $39\times$ \\
0.1   & 0.040 / 0.098 / 0.105 & 0.002 / 0.001 / 0.003 & $43\times$ \\
1.0   & 0.068 / 0.044 / 0.066 & 0.0002 / 0.0005 / 0.0004 & $161\times$ \\
\bottomrule
\end{tabular}
\caption{Per-seed spread of the single-model \IC\ versus AUC, per (arm, fraction) cell;
the ratio column is the mean-over-arms ratio. Training data grows $330\times$ down to the
last row: the AUC spread collapses ${\approx}30\times$, the \IC\ spread only
${\approx}3\times$.}
\label{tab:seednoise}
\end{table}

\paragraph{(b) The offset is weakly pinned by the loss, so training noise concentrates in
it.} Expanding the weighted cross-entropy about its minimum, a global logit shift
$\delta$ costs $\Delta L \approx \tfrac12\,\delta^2 \sum_i w_i c_{y_i}\, p_i(1-p_i)$: the
curvature is carried only by events near the decision boundary, and for a classifier at
AUC $0.83$ most events are confident, $p(1-p)\to 0$. Numerically, $\delta = 0.05$ (a $5\%$
\IC\ error, the size of the full-statistics seed scatter) costs
$\Delta L \sim 10^{-4}$ --- invisible to early stopping and to best-validation checkpoint
selection, so the frozen-in offset is effectively a per-seed random draw from the
training trajectory. Evaluation adds a further monotone jitter: EveNet predicts from an
exponential moving average of the \emph{parameters}, while normalisation-layer running
buffers are not averaged, an offset-and-gain perturbation by construction. The diagnostic
that separates this from sampling noise is the last row of Table~\ref{tab:seednoise}:
$330\times$ more training data shrinks the \IC\ seed spread only ${\approx}3\times$
(pure sampling would give $\sqrt{330}\approx 18\times$). The residual is optimisation and
selection stochasticity, which no single training removes --- and which ensembling
\emph{does} remove, because the offsets are approximately zero-centred across seeds
(Sec.~\ref{sec:closure}).

\paragraph{(c) The test side compounds it.} AUC is an average of a bounded comparison over
$N_0 N_1 \approx 1.7\times10^{9}$ pairs: test-statistical error ${\approx}0.002$. \IC\ is
the weighted mean of the heavy-tailed $\hat r$ over $4.9\times10^{4}$ reference events:
floor $0.012$--$0.019$. Per cell, the measured AUC effect ($0.015$) stands at
${\approx}7\sigma$ of its total noise; a plausible calibration effect ($0.03$) stands at
$<1\sigma$ of its. The same 90 runs power one axis and not the other by construction of
the metrics, not by bad luck.

\paragraph{(d) The regime matters: reweighting versus discrimination.} When the two
samples are nearly identical --- OmniFold-style data-versus-simulation
reweighting~\cite{Mikuni:2024omnilearn,H1:2025prelim}, or a nominal-versus-shifted
systematic variation --- the classifier sits at $p \approx 0.5$ everywhere, exactly where
the offset curvature of (b) is maximal: calibration is \emph{cheap} in that regime, and
the H1 measurement indeed reports training and initialisation variance negligible against
data statistics. Well-separated signal-versus-background, as here, is the opposite corner.
The published foundation-model evidence (validation loss, discrimination, distances
between normalised unfolded distributions) therefore does not --- and structurally cannot
--- speak to absolute ratio calibration at low $N$, which is the axis NSBI uniquely
needs. The corollary for this programme: the natural near-term calibration win for a
pre-trained model is the systematic-variation task of N3, whose ratios are near unity,
not the signal-extraction ratio.

\paragraph{(e) The arithmetic behind $K{=}20$.} At $K{=}5$ the seed term exceeds the test
floor by $4$--$7\times$ at reduced fractions. $K{=}20$ brings the per-cell error to
$0.02$--$0.06$, which resolves a consistent arm offset of ${\gtrsim}0.05$ at $2\sigma$
across the five fractions --- the scale suggested by the discrimination gains --- while
the seed count at which the seed term would meet the test floor is
$K^{*} = (\hat\sigma/\sigma_{\mathrm{test}})^2 \approx 40$--$400$: seeds beyond
${\approx}50$ buy little without a larger test split. Pairing the bootstrap draws across
arms removes the common test fluctuation from arm differences at zero cost. These three
numbers fix the (P5) design; Sec.~\ref{sec:p5} reports the executed campaign.

% =====================================================================
\section{Result 3 --- the (P5) verdict: no calibration gain at low statistics}
\label{sec:p5}

\paragraph{Execution.} Seeds $5$--$19$ were trained at fractions
$\{0.003, 0.01, 0.02, 0.05, 0.1\}$ (225 runs, identical data, recipe,
best-validation selection and test split; $\approx 13$ GPU-node-hours). One training
diverged --- finetune, fraction $0.05$, seed $12$: the validation loss went NaN near
epoch 12, \emph{displayed as $0.000$}, and early stopping retained the diverged state as
the best checkpoint; its predictions were $100\%$ non-finite. Members with non-finite
outputs or unreadable predictions are excluded by an objective rule applied identically
to all arms; three cells run at $K{=}19$ (one member lost in each of finetune-$0.05$,
frozen-$0.003$, scratch-$0.1$ --- one per arm, favouring none). All error propagation
uses each cell's actual $K$.

\paragraph{Verdict.} Applying (P5) exactly as registered
(Table~\ref{tab:p5verdict}): fine-tuning shows \emph{no} calibration data-efficiency
gain --- any effect is bounded at $|\Delta| \lesssim 0.05$, with a negative central
value --- and the frozen arm \emph{crosses the pre-registered harm threshold}: its
ensemble normalisation error at reduced fractions is worse than from-scratch by
$0.079 \pm 0.031$, negative in four of the five fractions. The secondary \SCr\ statistic
agrees in sign for both arms and is individually significant for frozen: the pre-trained
arms' ratio \emph{shapes} are also worse than scratch at low $N$
(Table~\ref{tab:p5cells}). An independent re-implementation of the verdict arithmetic
and a cross-check of every ensemble value against the per-seed closure outputs
(linearity, agreement $\sim 10^{-7}$) confirm the numbers.

\begin{table}[h]\centering\small
\begin{tabular}{@{}lccl|cc@{}}
\toprule
 & $\Delta_{|\IC|} \pm \sigma_\Delta$ & $z$ & verdict (as registered)
 & $\Delta_{\SCr} \pm \sigma$ & $z$ \\
\midrule
finetune & $-0.022 \pm 0.026$ & $-0.9$ & bounded: $|\Delta| \lesssim 0.05$
         & $-0.038 \pm 0.019$ & $-2.0$ \\
frozen   & $-0.079 \pm 0.031$ & $-2.6$ & \textbf{harms calibration}
         & $-0.069 \pm 0.021$ & $-3.2$ \\
\bottomrule
\end{tabular}
\caption{The (P5) statistics ($\Delta > 0$ would mean the pre-trained arm is better
calibrated than scratch). Primary $|\IC|$ column: paired-bootstrap $\oplus$ per-seed
terms, as registered. Secondary \SCr: unpaired bootstrap $\oplus$ single-model seed-spread
proxy at each cell's $K$; reported, not decisive.}
\label{tab:p5verdict}
\end{table}

\begin{table}[h]\centering\small
\begin{tabular}{@{}lccc|ccc@{}}
\toprule
 & \multicolumn{3}{c|}{$|\IC|$ (ensemble $\pm$ boot $\oplus$ seed$/\sqrt{K}$)}
 & \multicolumn{3}{c}{\SCr} \\
fraction & finetune & frozen & scratch & finetune & frozen & scratch \\
\midrule
0.003 & $0.022 \pm 0.066$ & $0.200 \pm 0.085$ & $0.073 \pm 0.061$ & 0.209 & 0.237 & 0.212 \\
0.01  & $0.169 \pm 0.042$ & $0.195 \pm 0.043$ & $0.001 \pm 0.047$ & 0.231 & 0.305 & 0.119 \\
0.02  & $0.111 \pm 0.024$ & $0.114 \pm 0.036$ & $0.110 \pm 0.037$ & 0.188 & 0.178 & 0.123 \\
0.05  & $0.050 \pm 0.029$ & $0.132 \pm 0.065$ & $0.035 \pm 0.027$ & 0.087 & 0.147 & 0.043 \\
0.1   & $0.109 \pm 0.024$ & $0.102 \pm 0.028$ & $0.130 \pm 0.032$ & 0.104 & 0.107 & 0.129 \\
\bottomrule
\end{tabular}
\caption{$K{=}20$ ensemble closure per cell. The ensemble AUC ordering on the same cells
is unchanged from Sec.~\ref{sec:auc}: pre-trained arms $+0.013$ to $+0.021$ above scratch
at every fraction $\ge 0.01$ (and $-0.009$ to $-0.014$ at the optimiser-transient
$0.003$ tier).}
\label{tab:p5cells}
\end{table}

\paragraph{Interpretation.} The two axes of Sec.~\ref{sec:power} now carry opposite
signs on the same networks and the same test events: pre-training improves event
\emph{ordering} at every usable fraction while degrading the absolute ratio
normalisation and shape at low $N$. The mechanism follows Sec.~\ref{sec:power}(b): the
calibration offset is weakly pinned by the loss, so the offset the \emph{initialisation}
carries survives small-$N$ training --- a from-scratch network starts unbiased and
settles near its own data's cross-entropy balance, while the pre-trained backbone imports
a normalisation prior from a different corpus that $2$k--$66$k events cannot repin.
Crucially this is an arm-level \emph{bias}, not zero-centred seed noise: ensembling ---
which resolved everything at $K{=}5$ --- cannot remove it, and it became visible
precisely when $K{=}20$ beat the seed noise below its size. It is a low-$N$ phenomenon:
at full statistics the ordering inverts and frozen is the best-calibrated configuration
(Sec.~\ref{sec:closure}).

\paragraph{Post-hoc annotations} (labelled as such; the verdict is mechanical):
(a)~the finetune average is dominated by the $0.01$ cell and the $0.1$ cell is positive
--- the finetune deficit is fraction-heterogeneous, the frozen deficit is not
(negative in $4/5$); (b)~the scratch $0.01$ value $|\IC| = 0.0007$ is a folded
$|\cdot|$ statistic sitting on zero with a seed error of $0.045$ --- its point value
should not be over-read, and the (P5) error propagation carries its full uncertainty;
(c)~at fraction $0.1$ both differences turn mildly positive, consistent with the
full-statistics inversion.

% =====================================================================
\section{Corrections applied during analysis}
\label{sec:corr}

Four hypotheses about the apparent calibration failure were tested. Three were our own
conversion errors; the fourth was refuted before it was believed. Recorded because the
first three each changed the numbers materially and the fourth would have cost 15 GPU jobs.

\paragraph{C1--C3: the ratio prior convention (a chain of three).} EveNet's classification
loss is
$\sum_i c_{y_i} w_i \mathrm{CE}_i / \sum_i c_{y_i} w_i$, with
$c = $ \texttt{normalization.pt["class\_balance"]} and $w$ the signed per-event weight
(\texttt{apply\_event\_weight} is true). Its population optimum is
$p/(1-p) = c_1\,\mathrm{d}W_1(x) / (c_0\,\mathrm{d}W_0(x))$, and since
$r = [\mathrm{d}W_1/W_1]/[\mathrm{d}W_0/W_0]$,
%
\begin{equation}
r(x) \;=\; \frac{p}{1-p}\cdot\frac{c_0 W_0}{c_1 W_1},
\qquad W_c = \texttt{class\_counts}[c] \;\;(\text{signed, train split}).
\label{eq:prior}
\end{equation}
%
We used, in order: $W_0/W_1$ from the test split (C1, wrong --- omits $c$ entirely);
$1.0$ (C2, ``class-balanced'' --- correct only if $c$ were inverse-frequency); and finally
Eq.~\eqref{eq:prior} (C3). The key fact is that \texttt{class\_balance} is \emph{inverse
effective number of samples}, not inverse frequency --- verified by reconstructing the
stored vector bit-for-bit in \texttt{float32} from \texttt{class\_counts} via
$\beta = 1-1/N$ effective-number reweighting --- so it cancels only 77.9\% of the class
imbalance in log space. Numerically
$f = (1.2616\times11351.2)/(0.7384\times22569.1) = \mathbf{0.8593}$, against an empirically
required $0.8702 \pm 0.0498$ measured over the 15 full-statistics runs
($0.22\sigma$; no arm dependence). C2 inflated $|\IC|$ by ${\approx}3\times$ and, because a
constant scale error shifts every bin equally, inflated \SCr\ by the same amount --- which
is why the first pass showed closure \emph{flat} in $N$ rather than converging.

\paragraph{C4: refuted --- the diffusion-noised objective.} We hypothesised that EveNet's
classification head, trained with \texttt{noise\_prob}$=1.0$ but evaluated at the clean
$t=0$ endpoint, is not a proper scoring rule there, and planned a 15-job ablation with
\texttt{noise\_prob}$=[0,0]$. Source audit refuted this before the result was believed:
\texttt{noise\_prob} is read at exactly one line, inside a \texttt{generation} branch that
three flags in our own configs disable (\texttt{ReconGeneration.include: false},
\texttt{generation-recon: [0,0]}, \texttt{classification-noised: [0,0]}). The branch that
ran, \texttt{deterministic}, sets \texttt{full\_time}$=0$ --- \emph{identical} to predict
time. There was no train/predict mismatch and the ablation would have been a no-op.

\paragraph{Other fixes.} Closure is evaluated with signed weights (the measure the loss
used); predict configs inherited an unresolvable relative \texttt{options.yaml} path and
lacked the head-inclusion block, so every prediction task initially died before loading a
checkpoint.

% =====================================================================
\section{Caveats}
\label{sec:caveats}

\begin{itemize}[leftmargin=2.2em]
\item \textbf{No sensitivity claim.} No $\sigma_{\kl}$ statement is made or implied; the
  stock schema discards $\tau$-identification and FastMTT $m_{\tau\tau}$ (equally in all
  arms), and the MC-statistics estimator question is open
  upstream~\cite{Shyamsundar:2025}.
\item \textbf{One channel, one hypothesis pair} (\kl$=1$ vs.\ $5$; $\mu\tau_h$, $e\tau_h$).
  The \kl$=0$ pair and the systematic-variation task are untouched.
\item \textbf{Fraction-1.0 cells remain $K{=}5$} (the extension covered reduced
  fractions only); full-statistics comparisons carry the original seed-noise floor.
\item \textbf{The (P5) verdict is recipe- and task-conditional}: it is measured for this
  fine-tuning recipe, this hypothesis pair, and this test-set size; a calibration-aware
  fine-tuning objective (e.g.\ re-pinning the output layer's normalisation before
  training) could change it --- that is a constructive question for the authors, not an
  assumption.
\item \textbf{Frozen passes by a thin margin} ($|\IC| = 0.0527$ against a gate of $0.0536$).
  It should not be over-read as a robust pass.
\item \textbf{Population mismatch vs.\ the ceiling only} (Sec.~\ref{sec:ceiling}); all
  arm-to-arm comparisons share identical events.
\item \textbf{The 0.003 tier} is optimiser-transient-dominated (Sec.~\ref{sec:auc}).
\end{itemize}

% =====================================================================
\section{Next steps}
\label{sec:next}

\begin{enumerate}[label=\textbf{N\arabic*}.,leftmargin=2.6em]
\item \textbf{[P1, S --- done]} Bootstrap the ensemble closure metrics: implemented as a
  Poisson($\mu{=}1$) test-set bootstrap with draws paired across arms
  (\texttt{eval\_ensemble.py -{}-bootstrap}, run automatically by the \texttt{eval}
  stage); the seed-extension campaign it feeds is pre-registered as (P5).
\item \textbf{[P1, S]} G1 corner-encoding A/B (\tauh\ at the vacant flag corner
  $(0,0,\pm1)$): the NPZs already exist for both encodings.
\item \textbf{[P1, M]} Money Plot 2 --- systematic-variation data efficiency, using the
  analytic injector where closure is exact by construction. This is the regime the FM's
  low-statistics claim was actually aimed at, and it is untested.
\item \textbf{[P2, S]} Low-fraction recipe ablation (batch size / learning rate at $0.003$)
  to convert the AUC inversion from artefact to measurement.
\item \textbf{[P2, S]} Re-measure the ceiling on the valid-$b$-pair population to remove the
  asterisk in Sec.~\ref{sec:ceiling}.
\end{enumerate}

\paragraph{For the EveNet authors.} (i)~The ratio conversion of Eq.~\eqref{eq:prior} is not
documented anywhere in the repository and is easy to get wrong in three distinct ways
(Sec.~\ref{sec:corr}); for any likelihood-ratio use of the classifier head this is the
single most valuable thing to document, or to expose as a utility. (ii)~Our G1 result will
be the first data on reusing \emph{vacant flag corners} for object types absent from
pre-training (\tauh, $\gamma$) --- would they endorse this as the zero-surgery interim?
(iii)~Would they consider retrofitting an explicit type-embedding table, initialised from
the current flag patterns so it is function-preserving on day one, making future object
types an append-a-row operation rather than schema surgery?
(iv)~Bug report from the (P5) campaign: a diverged training's NaN validation loss
displays as $0.000$, which early stopping and best-checkpoint selection then treat as
the best score --- the diverged state wins checkpoint selection silently. A NaN guard in
the trainer (terminate, and never checkpoint a non-finite loss) would fix it.
(v)~The Sec.~\ref{sec:p5} result suggests any FM claim aimed at likelihood-ratio use
should be powered and reported on the calibration axis, not only on discrimination
metrics --- and that the near-unity systematic-variation regime, not signal extraction,
is where pre-trained calibration gains are most plausible.

% =====================================================================
\begin{thebibliography}{9}\small

\bibitem{Hsu:2026evenet}
S.-C. Hsu et al., ``EveNet: A Foundation Model for Particle Collision Data Analysis,''
arXiv:2601.17126 [hep-ex] (2026).

\bibitem{Elsharkawy:2026}
Y. Elsharkawy, J. Birk, V. Mikuni, W. Bhimji, G. Kasieczka, B. Nachman,
``Pre-Training for Simulation-Based Science: A Study on Jet Foundation Model Training
Objectives,'' arXiv:2606.14870 (2026).

\bibitem{Mikuni:2024omnilearn}
V. Mikuni, B. Nachman, ``OmniLearn: A Method to Simultaneously Facilitate All Jet Physics
Tasks,'' arXiv:2404.16091 [hep-ph] (2024).

\bibitem{H1:2025prelim}
H1 Collaboration, ``Towards Unfolding All Particles in High $Q^2$ DIS Events,''
H1prelim-25-031, DIS 2025 (2025).

\bibitem{Shyamsundar:2025}
P. Shyamsundar et al., critique of MC-statistical uncertainty estimation in NSBI,
arXiv:2505.19156 (2025).

\end{thebibliography}

\end{document}
